\section{Large-sample theory for the local CDF}
\label{sec:theory}

For a fixed bandwidth, the inferential target is the kernel-localized
distribution \(F_{h,0}(\cdot\mid x_0)\); Section~\ref{sec:applications} uses
study-specific controlled-sampling benchmarks. We
establish that estimating the
borrowing weight has no first-order cost, derive a uniform CDF limit and a
multiplier approximation, and then characterize the efficiency of surrogate
borrowing. Replacing \(F_{h,0}\) by the point target
\(F_0(\cdot\mid x_0)\) requires localization bias to be negligible at the
stochastic scale.

\subsection{Conditions and local information}

Write
\[
N_{M,h}=Mp_Mh^d,
\qquad
\kappa_K=
\frac{\int K^2(u)\,du}{\{\int K(u)\,du\}^2}.
\]
The quantity \(N_{M,h}\) is the expected order of the verified sample in the
target neighborhood. It, rather than \(M\), determines the stochastic rate.

For a threshold-indexed function \(b\), define the local verification metric
\[
\|b\|_{\mathrm{ver}}^2
=
\sup_{y\in\mathcal Y}
\frac{p_Mh^d}{\mu_h^2}
E\left[
K_h^2(X-x_0)
\left\{\frac1{\pi_M(W)}-1\right\}b_y^2(W)
\right].
\]

\begin{highlevelassumption}[Observed-data identification and verification]
\label{ass:verification}
Within each row \(M\), the observations are independent and identically
distributed, the joint law of \((W,Y)\) is fixed, and \(R\perp Y\mid W\). The
known verification probability satisfies \(\pi_M(W)=p_Me_M(W)\), where
\(p_M=E\{\pi_M(W)\}\to p\in[0,1]\), \(E\{e_M(W)\}=1\), and
\(0<c_e\le e_M(W)\le C_e<\infty\) almost surely.
\end{highlevelassumption}

\begin{highlevelassumption}[Localization and the local-information rate]
\label{ass:localization}
The point \(x_0\) is interior to the support of \(X\), and \(f_X\) is positive
and continuous there. The nonnegative kernel is bounded, symmetric, compactly
supported, and of bounded variation, with positive finite first and second
integrals. The bandwidth satisfies \(h\to0\) and \(N_{M,h}\to\infty\).
\end{highlevelassumption}

\begin{highlevelassumption}[Candidate paths and sample separation]
\label{ass:candidates}
Regression fitting, weight selection, and CDF evaluation use disjoint samples
whose limiting proportions are positive. The fitted threshold regressions are
uniformly bounded and converge in the local verification metric to
deterministic triangular paths \(q_M\) and \(m_M\). The fold-specific oracle
paths converge in that metric to a common population-minimizing path \(a_M\),
and \(a_M\) converges locally and uniformly over thresholds in conditional
\(L_2\) to a bounded path \(a_\infty\).
\end{highlevelassumption}

\begin{highlevelassumption}[Transfer-path curvature and degeneracy]
\label{ass:transfer-path}
For each rotation, the normalized curvature
\(\mathcal H_{h,k}=p_MH_{M,h,k}\) converges to
\(H_{\infty,k}\in[0,\infty)\). The stabilizer satisfies
\(\varepsilon_{H,M}\downarrow0\) and
\(p_M\varepsilon_{H,M}=o(H_{\infty,k})\) when \(H_{\infty,k}>0\).
Write \(\mathcal H_{M,h}\) for the analogous normalized curvature of the
deterministic difference \(m_M-q_M\). Then
\[
\|m_M-q_M\|_{\mathrm{ver}}^2
\le \omega(\mathcal H_{M,h})+o(1),
\]
where \(\omega\) is nondecreasing, continuous at zero, and \(\omega(0)=0\).
The same relation holds in probability after replacing the deterministic
difference and curvature by \(\widehat m_k-\widehat q_k\) and
\(\mathcal H_{h,k}\), respectively.
\end{highlevelassumption}

\begin{highlevelassumption}[Threshold-process regularity]
\label{ass:process}
\begin{enumerate}[label=(\alph*),leftmargin=2em]
\item The threshold set is compact. Conditional CDF increments are locally
Lipschitz in the threshold, and augmentation-path increments are locally
H\"older in conditional \(L_2\). The path-oracle and fitted-difference classes
are uniformly VC-type, have square-integrable envelopes satisfying the
triangular-array Lindeberg condition, and have fitted-difference envelopes that
vanish in the local \(L_2\) metric.
\item The local covariance integrand converges uniformly to a continuous
kernel \(\Gamma_p\) defining a tight Gaussian process. For simultaneous bands,
the distribution of \(\|\mathbb G_p\|_{\mathcal Y}\) is continuous at the
target critical value.
\end{enumerate}
\end{highlevelassumption}

Here \(e_0(W)\) denotes the local limit of \(e_M(W)\) entering the covariance
integrand in Assumption~\ref{ass:process}(b).

\begin{highlevelassumption}[Point-target and quantile regularity]
\label{ass:point-quantile}
Uniformly over thresholds, the target CDF and covariate density have bounded
continuous second derivatives in \(x\) near \(x_0\). For quantile inference,
\(F_0(\cdot\mid x_0)\) is continuously differentiable and strictly increasing
on the threshold region, its values remain in an interior probability range,
and \(f_{Y\mid X}(\cdot\mid x_0)\) is bounded above and away from zero there.
\end{highlevelassumption}

Supplementary Section S1 gives the primitive conditions and their exact map to
Assumptions~\ref{ass:verification}--\ref{ass:point-quantile}. Write
\(\Delta_y(W)=a_{\infty,y}(W)-m_{0,y}(W)\).

\subsection{Adaptive estimation of the local CDF}

We first compare the feasible estimator with a path oracle. Let
\(\widehat F_h^{\mathrm{or}}\) denote the sample-split estimator that uses the
same fitted candidate regressions as the feasible estimator and replaces each
selected weight by its fold-specific population minimizer.

\begin{theorem}[Adaptive path-oracle equivalence]
\label{thm:oracle-equivalence}
Under Assumptions~\ref{ass:verification}--\ref{ass:transfer-path} and
Assumption~\ref{ass:process}(a), the criterion regret after multiplication by
\(p_M\) is \(O_{\mathbb P}(N_{M,h}^{-1/2}+\varepsilon_M)\), where
\(\varepsilon_M=p_M\varepsilon_{H,M}\). If the limiting curvature is positive,
\(\widehat\lambda_{h,k}-\lambda_{h,k}^{\mathrm{or}}=
O_{\mathbb P}(N_{M,h}^{-1/2})\). If, eventually,
\(\lambda_{h,k}^{\mathrm{or}}\in[\delta,1-\delta]\) for some fixed
\(\delta>0\), criterion regret is \(O_{\mathbb P}(N_{M,h}^{-1})\). If
curvature tends to zero, the
fitted correction vanishes in the local process metric. In both cases,
\[
\sqrt{N_{M,h}}
\left\|
\widehat F_h^{\mathrm{raw}}-\widehat F_h^{\mathrm{or}}
\right\|_{\mathcal Y}
=o_{\mathbb P}(1).
\]
\end{theorem}

The theorem separates adaptation from CDF estimation. The selection step is
first-order negligible whether the candidate path has positive curvature or
collapses to the anchor. The remaining stochastic rate depends on local
verified information even when the overall verification fraction decreases.
Define
\begin{align}
\mathbb Z_M(y)
&=
\sqrt{N_{M,h}}
\{\widehat F_h^{\mathrm{raw}}(y\mid x_0)-F_{h,0}(y\mid x_0)\},
\label{eq:scaled-cdf-process}\\
\mathbb Z_M
&\Rightarrow
\mathbb G_p
\quad\text{in }\ell^\infty(\mathcal Y).
\label{eq:process-limit}
\end{align}

\begin{theorem}[Uniform convergence of the kernel-localized CDF]
\label{thm:process}
Under Assumptions~\ref{ass:verification}--\ref{ass:process},
\eqref{eq:process-limit} holds for a tight
mean-zero Gaussian process with covariance kernel
\begin{align}
\Gamma_p(y,z)
&=
\frac{\kappa_K}{f_X(x_0)}
E\Bigg[
\frac{C_{yz}(W)}{e_0(W)}
+
\left\{\frac1{e_0(W)}-p\right\}
\Delta_y(W)\Delta_z(W)
\nonumber\\[-2pt]
&\hspace{24mm}
+p\{m_{0,y}(W)-F_0(y\mid x_0)\}
  \{m_{0,z}(W)-F_0(z\mid x_0)\}
\,\Bigm|\,X=x_0
\Bigg].
\label{eq:cov-limit}
\end{align}
\end{theorem}

The covariance identifies the part that adaptive borrowing can change: the
term involving \(\Delta_y\Delta_z\) is the cost of augmentation error. The
remaining terms come from verified-outcome variation and heterogeneity within
\(X=x_0\); the latter vanishes when \(p=0\). Scarcity enters through \(p_M\) in
the rate, while \(e_0^{-1}\) records where verification occurs locally. Equal
verified counts can therefore yield different precision at \(x_0\).

\subsection{Simultaneous bands and conditional quantiles}

Let \(\mathbb G_M^\xi\) be the evaluation-fold Rademacher multiplier process
formed from the fitted influence contributions. The Supplement gives its exact
construction and the estimated-process replacement argument.

\begin{theorem}[Multiplier process]
\label{thm:bootstrap}
Under Assumptions~\ref{ass:verification}--\ref{ass:process},
\begin{equation}
\mathbb G_M^\xi
\Rightarrow_{\mathbb P}
\mathbb G_p
\quad\text{in }\ell^\infty(\mathcal Y).
\label{eq:bootstrap-limit}
\end{equation}
If \(c_{1-\alpha}^\xi\) is the conditional \((1-\alpha)\)-quantile of the
multiplier supremum, then
\begin{equation}
\Pr\left[
F_{h,0}(y\mid x_0)
\in
\left\{
\widehat F_h^{\mathrm{raw}}(y\mid x_0)
\pm
\frac{c_{1-\alpha}^\xi}{\sqrt{N_{M,h}}}
\right\}
\text{ for every }y\in\mathcal Y
\right]
\to1-\alpha.
\label{eq:band}
\end{equation}
\end{theorem}

One multiplier critical value therefore gives simultaneous coverage over the
threshold region. This conclusion concerns \(F_{h,0}\). The next result states
the additional bias and shape conditions needed for the scientific point
target and its quantiles.

\begin{corollary}[Point-target CDF and conditional quantiles]
\label{cor:F0}
\label{cor:quantile}
Suppose Theorems~\ref{thm:process} and \ref{thm:bootstrap} and
Assumption~\ref{ass:point-quantile} hold. If
\(\sqrt{N_{M,h}}h^2\to0\), the process limit and confidence band remain valid
with \(F_0(\cdot\mid x_0)\) in place of \(F_{h,0}(\cdot\mid x_0)\). If the point
target is strictly increasing on the threshold region with conditional density
bounded above and away from zero, increasing rearrangement is first-order
neutral and, for compact \(\mathcal T\) in the target probability range,
\begin{align}
\sqrt{N_{M,h}}
\{\widehat Q_h(\cdot\mid x_0)-Q_0(\cdot\mid x_0)\}
&\Rightarrow
\mathbb Q_p,
\label{eq:quantile-limit}\\
\mathbb Q_p(\tau)
&=-
\frac{\mathbb G_p\{Q_0(\tau\mid x_0)\}}
{f_{Y\mid X}\{Q_0(\tau\mid x_0)\mid x_0\}}.
\label{eq:quantile-gaussian-limit}
\end{align}
\end{corollary}

The fixed-bandwidth simulations evaluate \(F_{h,0}\). In
Section~\ref{sec:applications}, EPA and VitalDB use the full-data empirical
analogue \(F_{h,\mathrm{emp}}\), whereas the synthetic AI study uses its
population target \(F_{h,0}\). The feasible point-estimation experiments compare with
\(F_0\).

\subsection{Efficiency of surrogate borrowing}

The canonical gradient identifies the efficiency benchmark, and the covariance
identity below gives the exact cost of augmentation error.

\begin{theorem}[Canonical gradient and local efficiency bound]
\label{thm:efficiency}
Fix \(M\) and \(h\). Under Assumption~\ref{ass:verification},
\(\mu_h>0\), and
\(E\{K_h^2(X-x_0)/\pi_M(W)\}<\infty\), the canonical gradient for the
regularized target at threshold \(y\) is
\begin{equation}
\phi_{M,h,y}^{\mathrm{eff}}(O)
=
\frac{K_h(X-x_0)}{\mu_h}
\left[
m_{0,y}(W)
+\frac{R}{\pi_M(W)}\{Z_y-m_{0,y}(W)\}
-F_{h,0}(y\mid x_0)
\right].
\label{eq:efficient-if-threshold}
\end{equation}
Its covariance over any finite threshold collection is the semiparametric
efficiency bound. The same bound applies when the verification law is
unrestricted because the gradient is orthogonal to its tangent space.
\end{theorem}

This theorem establishes the efficiency benchmark for any feasible
augmentation. The bound uses the true threshold regression \(m_0(W)\).
Orthogonality to the verification-law tangent space gives the same benchmark
when a designed sampling probability is treated as known.

For augmentation \(a\), let \(\Omega_{M,h}^{a}(y,z)\) denote the covariance of
the corresponding scaled influence functions.

\begin{corollary}[Covariance cost of augmentation error]
\label{cor:cov-order}
Under Assumption~\ref{ass:verification}, if the relevant second moments are
finite, then
\begin{equation}
\Omega_{M,h}^{a}(y,z)-\Omega_{M,h}^{m_0}(y,z)
=
\frac{h^d}{\mu_h^2}
E\left[
K_h^2(X-x_0)
\left\{\frac1{\pi_M(W)}-1\right\}
\delta_y^a(W)\delta_z^a(W)
\right].
\label{eq:psd}
\end{equation}
The matrix formed by this difference is positive semidefinite over every finite
threshold set. Hence the oracle augmentation \(m_0\) attains the efficiency
bound, and augmentation error accounts exactly for the covariance excess.
\end{corollary}

This identity turns surrogate quality into a covariance comparison at the
target neighborhood. The adaptive path restricts that comparison to
augmentations between the anchor and the fitted surrogate regression.

\begin{theorem}[Local transfer criterion]
\label{thm:transfer}
Under Assumption~\ref{ass:verification}, conditional on the candidate
regressions, suppose the weighted squares defining \(J_{M,h}(0)\) and
\(H_{M,h}\) are integrable. The integrated augmentation variance is the
quadratic criterion in Section~\ref{sec:method}. Its minimizer is
\(\lambda_h^{\mathrm{or}}\) when \(H_{M,h}>0\), and full borrowing improves on
the anchor exactly when \(2G_{M,h}-H_{M,h}>0\). The path-oracle improvement is
\begin{equation}
J_{M,h}(0)-J_{M,h}(\lambda_h^{\mathrm{or}})
=
\begin{cases}
0, & H_{M,h}=0\text{ or }G_{M,h}\le0,\\[2pt]
G_{M,h}^2/H_{M,h}, & 0<G_{M,h}<H_{M,h},\\[2pt]
2G_{M,h}-H_{M,h}, & G_{M,h}\ge H_{M,h}>0.
\end{cases}
\label{eq:oracle-gain}
\end{equation}
\end{theorem}

The criterion compares the anchor and Full along the fitted correction path.
Its sign diagnoses full borrowing, while \(\lambda_h^{\mathrm{or}}\) chooses the
best point on that path. Neither statement implies finite-sample dominance over
the anchor, a distinction examined in the numerical studies.

The Supplement gives complete proofs, the finite-threshold specialization, the
limit of the efficiency bound, estimated-probability conditions, cluster
inference, and local-linear extensions.
